% TOP_PROBLEM: {"id": 30006804, "title": "Kawamata's K-Equivalence Implies D-Equivalence Conjecture", "category_id": 5, "category": {"id": 5, "name": "algebraic_geometry", "display_name": "Algebraic Geometry", "description": "Geometric objects defined by polynomial equations.", "slug": "algebraic-geometry", "order_index": 5, "created_at": "2026-07-31T15:26:25.671Z"}, "status": "open"}
% FIRST_POSTED: 2026-10-01
% ATTEMPT_STATUS: unresolved
% !TeX program = lualatex
% Definition and sources only; this file is not an LLM solution attempt.
\documentclass[11pt]{article}
\usepackage[margin=25mm]{geometry}
\usepackage{fontspec}
\setmainfont{Latin Modern Roman}
\usepackage{amsmath,amssymb,mathtools}
\usepackage{unicode-math}
\setmathfont{Latin Modern Math}
\usepackage{xurl,hyperref}
\hypersetup{colorlinks=true,urlcolor=blue}
\setcounter{secnumdepth}{0}
\setlength{\parindent}{0pt}
\setlength{\parskip}{5pt}
\setlength{\emergencystretch}{3em}
\begin{document}
\section*{\texorpdfstring{Kawamata's K-Equivalence Implies D-Equivalence Conjecture}{Kawamata's K-Equivalence Implies D-Equivalence Conjecture}}\label{30006804}
\subsection{Definitions and mathematical statement}
For smooth projective complex varieties $X,Y$, the selected $K$-equivalence means there is a common smooth $Z$ with projective birational morphisms $f:Z\to X$, $g:Z\to Y$ such that
$f^*K_X\sim g^*K_Y$, where $\sim$ is linear equivalence of divisors. The bounded derived category $D^b\operatorname{Coh}(X)$ is formed from bounded complexes of coherent sheaves by inverting quasi-isomorphisms. The assertion is
\[
 X\text{ and }Y\text{ are }K\text{-equivalent}
 \quad\Longrightarrow\quad D^b\operatorname{Coh}(X)\simeq D^b\operatorname{Coh}(Y)
\]
by an exact equivalence of triangulated categories. Only this implication is requested; replacing linear equivalence with numerical equivalence or asserting a converse would alter the target.
\par\smallskip\textit{Formulation and concept references:} \href{https://arxiv.org/abs/math/0205287}{[S1]}.

\subsection{Short English statement}
Do smooth projective varieties whose canonical divisors agree after birational pullback necessarily have equivalent derived categories of coherent sheaves?
\subsection{Research attempt: the curve case and tests for a correspondence functor}
\textbf{Status.} We prove the implication for smooth projective curves, compute an elementary birational case excluded by $K$-equivalence, and derive necessary categorical tests for a proposed functor. The general implication remains open in this attempt. The primary formulation and derived-category context are those of Kawamata [S1].

\subsubsection{1. Smooth projective curves}
A birational map $\varphi:X\dashrightarrow Y$ between smooth projective integral curves extends across every missing closed point of $X$. Indeed, its local ring is a discrete valuation ring, and the generic-point map into proper $Y$ extends uniquely by the valuative criterion of properness. These extensions agree with the rational map and glue, since $Y$ is separated. The birational inverse extends in the same way. Their compositions agree with the identity on dense open sets and therefore everywhere, again by separatedness. Consequently $X\cong Y$.

In the target's curve case, the common birational model already gives such a birational map, so this argument applies without needing further canonical-divisor information. Pullback along the resulting isomorphism is an exact equivalence of coherent sheaves; applying it term by term gives an exact equivalence of bounded derived categories. This proves the selected implication in dimension one. Dimension zero is immediate for integral complex varieties.

\subsubsection{2. Why an arbitrary blowup is not a test case for the conjecture}
Let $f:Z\to X$ be the blowup of a smooth center of codimension $c\ge2$, with exceptional divisor $E$. The standard local Jacobian calculation gives
\[K_Z=f^*K_X+(c-1)E.\]
For a line $\ell$ in a fiber $\mathbb P^{c-1}$, one has $E\cdot\ell=-1$, hence
\[(K_Z-f^*K_X)\cdot\ell=-(c-1)\ne0.\]
Thus these pullbacks are not even numerically equivalent on the natural common model, much less linearly equivalent. The difference also cannot become linearly trivial merely by further resolving this same birational correspondence: a proper birational morphism between smooth varieties has injective pullback on line bundles, since pushing forward its structure sheaf gives the structure sheaf. Thus no such refinement removes the discrepancy. This proves that the specified blowup correspondence does not exhibit a $K$-equivalence; no different birational identification of these varieties is being analyzed in this calculation.

This computation blocks an invalid shortcut: birational maps alone do not satisfy the premise. In codimension one the blowup of a Cartier center is an isomorphism, consistent with the zero discrepancy; it is not an exceptional new example.

\subsubsection{3. Necessary tests for a proposed equivalence}
For smooth projective $X$, define the Euler pairing on objects by
\[\chi_X(A,B)=\sum_i(-1)^i\dim_{\mathbb C}\operatorname{Hom}(A,B[i]).\]
An exact equivalence preserves each Hom space and shifts, hence preserves this pairing and induces an isomorphism on the Grothendieck group. Moreover the Serre functor is characterized by natural dualities $\operatorname{Hom}(A,B)^*\cong\operatorname{Hom}(B,S_XA)$. Conjugating this characterization by an equivalence and using uniqueness of a representing functor shows that an equivalence must intertwine $S_X$ and $S_Y$. In the smooth projective case these are $(-)\otimes\omega_X[\dim X]$ and its counterpart on $Y$.

These facts explain why canonical divisors are relevant, but they are necessary tests, not sufficient tests. Given a common model, the candidate $Rg_*Lf^*$ is a well-defined exact functor on the bounded coherent categories here. Equality $f^*K_X\sim g^*K_Y$ alone does not establish that its unit and counit are isomorphisms, nor that it is the correct kernel in every example. A proof must actually produce full faithfulness and essential surjectivity, possibly using a different correspondence.

\subsubsection{4. Audit and gap}
The curve argument uses properness to extend and separatedness to identify maps, so it does not extend unchanged to higher-dimensional indeterminacy loci. The blowup obstruction uses $c\ge2$ and the sign of $E$ on fiber lines. The categorical calculations preserve linear canonical data rather than replacing the premise by numerical equality. No construction here handles arbitrary higher-dimensional $K$-equivalent varieties; that is the exact remaining scope.

\subsection{Sources}
\begin{itemize}
\item[S1]Yujiro Kawamata, “D-Equivalence and K-Equivalence,” Journal of Differential Geometry 61 (2002), arXiv:math/0205287v3.
\url{https://arxiv.org/abs/math/0205287}.
\textit{Formulation source.}
\end{itemize}
\end{document}
